\begin{figure}
    \includegraphics[width=\linewidth]{Figures/Z.pdf}
    \caption{Schematic overview of the dataset generation pipeline. The diagram outlines the steps for constructing synthetic catalogues tailored to either photometric or spectroscopic selection criteria, using mock galaxies from \texttt{OpenUniverse2024} and \texttt{CosmoDC2}.}
    \label{fig:2}
\end{figure}

In the following, we outline our approach to utilising these advanced mock galaxy catalogues to produce several synthetic datasets. Each dataset is designed to replicate the photometric and spectroscopic selection functions of galaxy surveys for LSST photometric redshift analyses and to account for variations in selection effects. The datasets we generated are defined as follows:

\begin{itemize}
    \item \texttt{Application}: a general LSST-like photometric sample.
    \item \texttt{Degradation}: a compilation of multiple spectroscopic datasets derived from the \texttt{Application} datasets, used as part of the training sample.
    \item \texttt{Augmentation}: training datasets generated through simulations and added to the spectroscopic samples.
    \item \texttt{Combination}: complete training datasets produced through the integration of \texttt{Degradation} and \texttt{Augmentation} samples.
\end{itemize}

Figure~\ref{fig:2} presents a schematic overview of the pipeline, illustrating the principal stages from catalogue ingestion to the production of analysis-ready datasets. In order to incorporate different systematic effects into later modelling stages, we conduct \(501\) experiments that produce the four catalogues. The study involves a baseline experiment for machine learning training and \(N = 500\) supplementary experiments for statistical analysis, distinguished by differences in photometric noise realisations, spectroscopic selection functions, and samples of augmented galaxies, as outlined subsequently.

\subsection{\texttt{Application} datasets}

As mentioned earlier, the \texttt{OpenUniverse2024} catalogues serve as the foundation for replicating the photometric \texttt{Gold} selection as outlined in the LSST DESC Science Requirements Document \citep[SRD;][]{2018arXiv180901669T}. We refer to the outcome of our \texttt{Gold} selection as the \texttt{Application} datasets. We then apply an unsupervised machine‐learning technique for dimensionality reduction, the Self‐Organising Map (SOM), to map galaxy broadband SEDs onto a two‐dimensional grid, thereby facilitating subsequent analyses of spectroscopic degradation and data augmentation.

\subsubsection{\texttt{Gold} selection}

\begin{table*}
  \centering
  \begin{threeparttable}
    \caption{Statistical properties of the baseline \texttt{Application} dataset.}
    \label{tab:1}
    \begin{tabular}{lccccc}
      \hline
      Scenario & Number of galaxies & Number density & Effective number density\tnote{a} & Limiting magnitude\tnote{b} & Gold magnitude\tnote{c}  \\
      \hline
      LSST-Y1 & 2924005 & \(14.77 \, \mathrm{arcmin}^{-2}\) & \(11.14 \, \mathrm{arcmin}^{-2}\) & \(25.26 \) & \(24.05 \) \\
      LSST-Y10 & 7223012 & \(36.48 \, \mathrm{arcmin}^{-2}\) & \(27.95 \, \mathrm{arcmin}^{-2}\) & \(26.51 \) & \(25.30 \) \\
      \hline
    \end{tabular}
    \begin{tablenotes}
      \item[a] Effective number density is defined in \cite{2013MNRAS.434.2121C}.
      \item[b] \(i-\)band \(5\sigma\) point-source limiting magnitude under the prescribed observing conditions.
      \item[c] \texttt{Gold} magnitude: the \(i-\)band magnitude threshold used for \texttt{Gold} selection, as specified by LSST DESC SRD.
    \end{tablenotes}
  \end{threeparttable}
\end{table*}

Utilising the \texttt{OpenUniverse2024} catalogues, our first step is to compute realistic photometric uncertainties for observations at the LSST Y1 and Y10 depths. These uncertainties are obtained using the publicly available tool \citep[\href{https://github.com/jfcrenshaw/photerr}{\faGithub\,\texttt{PhotErr}};][]{2024AJ....168...80C}, which adopts the single-visit magnitude limits and expected numbers of visits per band from the survey specifications presented in \citet{2019ApJ...873..111I}. In our implementation, these quantities are used solely to derive uniform-depth Y1 and Y10 photometric-error models; we do not model LSST footprint variations or spatially varying depth. Only the Wide–Fast–Deep (WFD) component of the survey is considered, and the Deep Drilling Fields (DDFs) are not included. The resulting error prescriptions are applied to both point sources and extended sources using the galaxy morphological parameters described by \citet{2020A&A...642A.200V}.

In the \texttt{Application} datasets, each of the \(501\) experiments incorporates a separate random realisation of photometric noise, achieved through a unique seed to disturb the galaxy photometry. This method guarantees the creation of statistically independent synthetic datasets that mirror observational uncertainties. When calculating the perturbed magnitudes, we manually set the fluxes at the \(1 \sigma\) noise threshold for galaxies dimmer than the specified limiting magnitudes in each band, considering them as non-detections. This setting ensures that the distribution of galaxy colours remains numerically continuous, even when certain bands are not detected. 

One of the key aims of the \texttt{Gold} selection is to pinpoint galaxies with trustworthy detections. Consequently, the \texttt{Application} datasets developed subsequent to the \texttt{Gold} selection embody pragmatic photometric selection standards and serve as the groundwork for subsequent photometric-redshift modelling. Therefore, in order to eliminate very faint extended sources, we require the Signal‐to‐Noise Ratio (SNR) obtained from the combined \(r\)- and \(i\)-band photometry, (\(\mu \equiv \mathrm{SNR}_{i + r}\)), to meet the condition \(\mu > 10\). Additionally, we impose a restriction on the squared ratio of the galaxy size to the coadded point-spread function (PSF), expressed as \(\eta \equiv \left( R^{i + r} \,/\, R_{\mathrm{PSF}}^{i + r} \right)^2\), requiring \(\eta > 0.1\) in order to reduce contamination from unresolved point sources during practical analysis.

In addition, we exclude objects brighter than \(16 \) in \(i-\)band to mimic the masking of bright, saturated sources, and remove galaxies fainter than \(24.05\) for Y1 and \(25.30\) for Y10, in accordance with the \texttt{Gold} selection defined in the LSST DESC SRD. Additionally, to capture spatial sample variance, we randomly select half of the HEALPix sky pixels from \texttt{OpenUniverse2024} in each experiment.

Table~\ref{tab:1} provides a concise overview of essential information about the baseline \texttt{Application} datasets. Differences in the supplementary experiments are minimal and thus excluded for the sake of brevity.

\subsubsection{Self-Organising Map}

\begin{figure*}
    \centering
    \begin{subfigure}[b]{0.48\linewidth}
        \centering
        \includegraphics[width=\linewidth]{Figures/Y1/MAP.pdf}
        \caption{LSST Y1}
    \end{subfigure}
    \hspace{0.0\linewidth}
    \begin{subfigure}[b]{0.48\linewidth}
        \centering
        \includegraphics[width=\linewidth]{Figures/Y10/MAP.pdf}
        \caption{LSST Y10}
    \end{subfigure}
    \caption{Colour maps showing the average true redshift per SOM cell for the baseline experiment. The panels on the left correspond to SOMs trained on the LSST~Y1 photometric sample, and those on the right correspond to SOMs trained on the LSST~Y10 sample. Within each panel, the \texttt{Application}, \texttt{Degradation}, \texttt{Augmentation}, and \texttt{Combination} datasets occupy the upper-left, upper-right, lower-left, and lower-right quadrants, respectively. The \texttt{Degradation} dataset spans a noticeably smaller region of the SOM plane, reflecting the reduction in colour-magnitude coverage imposed by spectroscopic-selection effects. Meanwhile, although the \texttt{Augmentation} dataset broadly follows the structure of the \texttt{Application} sample, it departs more markedly in the Y10 case, where the underlying SED-model differences between \texttt{OpenUniverse2024} and \texttt{CosmoDC2} become increasingly evident.}
    \label{fig:3}
\end{figure*}

We employ the baseline \texttt{Application} dataset to train the SOM algorithm implemented via \texttt{Somoclu} \citep{2013arXiv1305.1422W} within the \texttt{RAIL} platform. This training compresses the high-dimensional space of a suite of representative galaxy broadband SEDs into a discrete two-dimensional map, with each cell corresponding to a sub-volume of the SED space. We choose the \(i-\)band as the reference, combining its apparent magnitudes with the colours derived relative to it. The SOM grids are configured as \(140 \times 140\) for Y1 and \(180 \times 180\) for Y10. These choices reflect a balance between achieving sufficient resolution to capture the mapping between galaxy broadband SEDs and redshift, and the practical limitations imposed by available computational resources. Although the present grid sizes are adequate for our analysis, future applications to petabyte-scale LSST datasets may require further optimisation, either through improved parallelisation and accelerated SOM training, or through carefully designed downsampling strategies for the observational catalogues.

Additionally, we adopt a Gaussian kernel with a standard deviation coefficient of \(0.5\) and a toroidal topology with hexagonal neurons for the training. The training is initialised using Principal Component Analysis (PCA) and then iterated for \(100\) epochs with linear cooling of both the neighbourhood radius and learning rate. The pre-trained SOM is subsequently applied to the remaining \(500\) statistical \texttt{Application} datasets, computing the Best Matching Unit (BMU) for each galaxy. To minimise computational costs, we refrain from retraining for each iteration and verify that the random photometric error realisation does not substantially affect the SOM mapping.

In each panel of Figures~\ref{fig:3}–\ref{fig:6}, the left panel corresponds to the Y1 scenario and the right to Y10, both illustrating the baseline experiment as a representative example of the full set of 501 realisations. In Figure~\ref{fig:3}, the upper‐left subplot shows the SOM map of mean true redshifts derived from the baseline \texttt{Application} dataset. Figures~\ref{fig:4} displays the \(u - g\) versus \(g - r\) and \(i - z\) versus \(z - y\) colour–colour diagrams, again with the baseline \texttt{Application} data in the upper‐left quadrants of each panel. Figure~\ref{fig:5} illustrates the distribution of \(g - z\) colour against the \(i\)-band magnitude, whereas Figure~\ref{fig:6} shows the galaxy distributions in relation to the true redshifts \(z_{\mathrm{true}}\). Note that the Y10 dataset contains a higher fraction of galaxies at \(z_\mathrm{true} \gtrsim 1.0\) and fainter magnitudes (\(i > 23 \)) than Y1, reflecting its deeper survey depth.  

\subsection{\texttt{Degradation} datasets}

\begin{figure*}
    \centering
    \begin{subfigure}[b]{0.48\linewidth}
        \centering
        \includegraphics[width=\linewidth]{Figures/Y1/COLOUR.pdf}
        \caption{LSST Y1}
    \end{subfigure}
    \hspace{0.0\linewidth}
    \begin{subfigure}[b]{0.48\linewidth}
        \centering
        \includegraphics[width=\linewidth]{Figures/Y10/COLOUR.pdf}
        \caption{LSST Y10}
    \end{subfigure}
    \caption{Colour--colour diagrams for the baseline experiment, shown for LSST~Y1 (left) and LSST~Y10 (right). The upper panels display \(u - g\) versus \(g - r\), while the lower panels show \(i - z\) versus \(z - y\). Each panel is subdivided into four quadrants corresponding to the \texttt{Application} (upper left), \texttt{Degradation} (upper right), \texttt{Augmentation} (lower left), and \texttt{Combination} (lower right) datasets. A logarithmic scale highlights structures in low-density regions of colour-magnitude space while maintaining the qualitative sense of relative densities.}
    \label{fig:4}
\end{figure*}

After constructing the \texttt{Application} datasets, we perform \(501\) realisations of a two-stage degradation procedure to mimic the selection effects that arise when compiling multiple spectroscopic surveys. This yields a suite of synthetic spectroscopic datasets representative of those plausibly accessible for LSST analysis. In the initial stage, we apply cuts in magnitude, redshift, and colour, and incorporate probabilistic sampling based on spectroscopic success rates. This produces \(501\) unique realisations of individual spectroscopic surveys, each reflecting the characteristics of different existing or planned spectroscopic programmes. The framework is flexible and can readily accommodate a wide variety of prospective spectroscopic strategies, rather than being tied to any single survey configuration.

In the second stage, we randomly sample and combine subsets of these realisations to generate \(501\) \texttt{Degradation} datasets, thereby emulating the compilation of multiple spectroscopic catalogues into a more representative training set for photometric-redshift modelling. The resulting \texttt{Degradation} datasets capture realistic variations in spectroscopic depth, colour-dependent selection biases, and redshift distributions, each of which is forward-modelled as a source of systematic uncertainty in the subsequent photometric-redshift estimation. However, our simplified prescriptions are not expected to fully capture all aspects of real spectroscopic compilations, but instead offer a forecasting tool for prospective LSST analyses.

\subsubsection{Spectroscopic selection}

For each realisation of the \texttt{Application} datasets, we draw an upper limit on the observed \(i\)-band magnitude, \(i_\ast\), uniformly from \(20  < i_\ast < 24 \) for Y1 and \(20  < i_\ast < 25 \) for Y10. These intervals mirror the depths of spectroscopic surveys likely to accompany LSST in its early and later operations.  Likewise, we assign an upper redshift cut, \(z_\ast\), by sampling uniformly within \(0.5 < z_\ast < 2.5\) for Y1 and \(0.5 < z_\ast < 3.0\) for Y10. 

This parameter configuration anticipates the increasing focus on faint and high-redshift galaxies in upcoming spectroscopic programmes. These encompass the accumulation of spectroscopy of galaxies through ongoing initiatives, as well as space-based missions like the James Webb Space Telescope \citep[JWST\footnote{\url{https://webbtelescope.org/home}};][]{2006SSRv..123..485G}, next-generation ground-based facilities -- such as the Extremely Large Telescope \citep[ELT\footnote{\url{https://elt.eso.org}};][]{2023ConPh..64...47P} and Giant Magellan Telescope \citep[GMT\footnote{\url{https://giantmagellan.org}};][]{2024SPIE13094E..17B} -- as well as preliminary Stage-V spectroscopic surveys, including DESI-II \citep{2022arXiv220903585S}, the MUltiplexed Survey Telescope (MUST\footnote{\url{https://must.astro.tsinghua.edu.cn}}) \citep{2024arXiv241107970Z}, and other potential candidates \citep{2022arXiv220904322S, 2025arXiv250307923B}.

Furthermore, deep-field spectroscopic programs frequently focus on optically red galaxies due to their robust redshift determination obtained from distinct spectral features, together with their more predictable selection functions, stronger large-scale clustering, and their detectability at greater depths. To simulate this preference, a magnitude-dependent cut in the \(g - z\) colour is applied, characterised by a linear relation with the \(i\)-band magnitude:
\begin{equation} \label{eqa:1} 
    i - i_\ast < \tan{\varphi_\ast} \left[ \left( g - z \right) - \left( g - z \right)_\ast \right].
\end{equation} 
Here, \(i_\ast\) denotes the previously defined upper limit in the \(i\)-band, \((g - z)_\ast\) the colour intercept and \(\varphi_\ast\) the slope angle. For both the LSST Y1 and Y10 scenarios, \((g - z)_\ast\) is drawn uniformly from \([1.0,\,3.0]\) and \(\varphi_\ast\) from \([0,\,\pi/2]\). 

This functional form is motivated by the colour-magnitude pre-selection strategies widely used in deep spectroscopic surveys such as DEEP2 \citep{2013ApJS..208....5N}, VVDS \citep{2005A&A...439..845L}, LEGA-C \citep{2021ApJS..256...44V}, and VUDS \citep{2015A&A...576A..79L}. These surveys often employ linear or piecewise-linear colour-magnitude boundaries to preferentially target optically red galaxies with strong continuum and absorption features, enabling robust redshift determination at faint magnitudes. Our parametrisation in Eq.~\ref{eqa:1} is therefore a simplified but flexible representation of these empirically motivated selection strategies.

In contrast, wide-field campaigns like \textit{Euclid} and DESI also target bluer star-forming galaxies that exhibit prominent emission lines, particularly at high redshift. This selection is mainly influenced by detectability criteria regarding emission line flux and redshift window instead of colour \citep{2024A&A...689A.166C, 2023AJ....165..126R}, making them less affected by a red-galaxy exclusion. In order to simulate the influence of more intricate selection functions within the multi-dimensional colour space, we randomly choose SOM cells at a rate \(r_\ast\), which is uniformly distributed across the interval \(r \in \left[0.2, 0.8\right]\). This is applied to both the LSST Y1 and Y10 scenarios, with galaxies in the unselected cells being disregarded. Although our probabilistic prior within this parametrised framework is deliberately simplistic, it provides useful flexibility to capture variations in colour thresholds across different surveys. More sophisticated, survey‐specific colour boundaries could be incorporated in future developments.

Moreover, we account for the impact of the spectroscopic success rate \(\lambda (i)\) by modelling it as a logistic function of the galaxy \(i-\)band apparent magnitude: 
\begin{equation} \label{eqa:2}
    \lambda (i) = \frac{1}{1 + \lambda_\ast \exp{\left[ i - i_\ast \right]}}.
\end{equation} 
In this context, \(i_\ast\) represents the upper threshold previously defined in the \(i\)-band magnitude for each \texttt{Application} dataset. To introduce additional stochasticity, we draw the parameter \(\lambda_\ast\) uniformly from the range \(0.5 < \lambda_\ast < 5.0\) for both Y1 and Y10 scenarios, thereby capturing different levels of spectroscopic survey completeness. 

The logistic form in Eq.~\ref{eqa:2} produces an almost complete success rate for the brightest sources (\(i \lesssim 18 \)) and a steep decline as the SNR decreases. This behaviour reflects the well-established magnitude/SNR dependence of spectroscopic completeness observed in surveys such as zCOSMOS \citep{2009ApJS..184..218L}, VIPERS \citep{2014A&A...566A.108G}, and DESI \citep{2023AJ....165..126R}, all of which report near-unity success rates at bright magnitudes followed by a rapid fall-off beyond their survey-specific limits. Our parametrisation is intended to capture this characteristic transition while remaining survey-agnostic and flexible enough to accommodate variations in depth and targeting efficiency across realisations.

Each realisation of the first‐stage spectroscopic dataset is obtained by sampling according to Equation~\ref{eqa:2}, yielding \(100000\) to \(200000\) galaxies for Y1 and \(250000\) to \(500000\) galaxies for Y10, in line with the expected expansion of spectroscopic coverage. These sample sizes are motivated by deep‐field spectroscopic surveys, which cover relatively small areas of sky. By contrast, wide‐field campaigns observe orders of magnitude more galaxies; including them in full would cause the spectroscopic training sample to be dominated by their shallower selection functions, introducing a pronounced imbalance between bright and faint galaxies in photometric‐redshift training. Consequently, we select a subset of the wide‐field spectroscopic datasets of comparable size to our deep‐field samples, as reflected in our configuration. Moreover, given finite computational resources, much larger spectroscopic samples would substantially increase the cost of photometric‐redshift training across multiple realisations.

\subsubsection{Compilation of multiple datasets}

Finally, every \texttt{Degradation} dataset is constructed by merging \(M\) randomly chosen realisations, where \(M\) is drawn from \([4,\, 6]\) for Y1 and \([7,\, 12]\) for Y10, and then downsampling to the same total galaxy counts as the corresponding experiment of spectroscopically-selected galaxies, thereby minimising computational overhead. This procedure produces \texttt{Degradation} datasets that encapsulate multiple realistic spectroscopic selection functions while preserving sufficient variation across experiments.  

The \texttt{Degradation} datasets for pre‐trained SOM cells are shown in the upper right subplots of each panel in Figure~\ref{fig:3}, illustrating how spectroscopic selection sculpts the galaxy distribution in SOM space. Notably, many cells, particularly those at higher redshift, become empty due to these selection effects. This sparsity is also evident in the colour–colour diagrams (upper right quadrants of each panel in Figure~\ref{fig:4}), in the magnitude and colour cuts depicted in Figure~\ref{fig:5}, and in the lack of galaxies with \(z_\mathrm{true} \gtrsim 1.5\) (dark‐blue solid lines) in Figure~\ref{fig:6}. If left uncorrected, such incomplete sampling can heavily bias subsequent photometric‐redshift estimates.

\subsection{\texttt{Augmentation} datasets}

To address selection effects introduced by the degradation process, we adopt and extend an augmentation strategy originally proposed by \cite{2024ApJ...967L...6M}, adapting it for use with a SOM. Our method comprises a dual-phase procedure to construct the \texttt{Augmentation} datasets. Initially, galaxies are sampled from the \texttt{CosmoDC2} catalogues to ensure their distribution in SOM space corresponds with that of the \texttt{Application} datasets, thus aligning their overall distributions within the multi-dimensional colour-magnitude space. In the second phase, we perform targeted augmentation to populate SOM cells that are under-represented or entirely unoccupied in the corresponding \texttt{Degradation} datasets.

\subsubsection{Weighted sampling}

\begin{figure*}
    \centering
    \begin{subfigure}[b]{0.48\linewidth}
        \centering
        \includegraphics[width=\linewidth]{Figures/Y1/DIAGRAM.pdf}
        \caption{LSST Y1}
    \end{subfigure}
    \hspace{0.0\linewidth}
    \begin{subfigure}[b]{0.48\linewidth}
        \centering
        \includegraphics[width=\linewidth]{Figures/Y10/DIAGRAM.pdf}
        \caption{LSST Y10}
    \end{subfigure}
    \caption{\(g - z\) colours versus \(i\)-band apparent magnitudes for various datasets are displayed for Y1 (left panel) and Y10 (right panel), under the baseline experiment. In each panel, the upper left, upper right, lower left, and lower right quadrants correspond to the \texttt{Application}, \texttt{Degradation}, \texttt{Augmentation}, and \texttt{Combination} datasets, respectively. A logarithmic scale enhances features in low-density areas of colour-magnitude space while preserving the perceived relative densities.}
    \label{fig:5}
\end{figure*}

Initially, \texttt{CosmoDC2} galaxies are processed using the same pipeline as for the \texttt{Application} datasets. This includes applying LSST photometric errors, perturbing apparent magnitudes, enforcing the LSST DESC SRD \texttt{Gold} selection, and randomly sampling over HEALPix sky pixels to account for spatial variations, under both Y1 and Y10 observing conditions.

Then, the pre-trained SOM is used to map each selected \texttt{CosmoDC2} galaxy to its BMU in SOM space. We calculate cell-wise occupancy for both the \texttt{Application} datasets and the selected \texttt{CosmoDC2} galaxies. We then assign sampling weights to the \texttt{CosmoDC2} galaxies given by the ratio of cell occupancies, ensuring that the synthetic samples reproduce the structure of the \texttt{Application} datasets in SOM space. This selection and re-weighting procedure is repeated over \(501\) independent experiments, one baseline and \(500\) statistical runs, each matched in size to their respective \texttt{Application} datasets to minimise excessive manipulation.

\subsubsection{Adaptive data augmentation}

\begin{figure*}
    \centering
    \begin{subfigure}[b]{0.48\linewidth}
        \centering
        \includegraphics[width=\linewidth]{Figures/Y1/HISTOGRAM.pdf}
        \caption{LSST Y1}
    \end{subfigure}
    \hspace{0.0\linewidth}
    \begin{subfigure}[b]{0.48\linewidth}
        \centering
        \includegraphics[width=\linewidth]{Figures/Y10/HISTOGRAM.pdf}
        \caption{LSST Y10}
    \end{subfigure}
    \caption{Normalised ensemble distributions of galaxy true redshifts for various datasets are displayed for for Y1 (left panel) and Y10 (right panel), under the baseline experiment. In each panel, the black, dark-blue, dark-red, dark-orange correspond to the \texttt{Application}, \texttt{Degradation}, \texttt{Augmentation}, and \texttt{Combination} datasets, respectively.}
    \label{fig:6}
\end{figure*}

To mitigate incomplete sampling within the SOM space of the \texttt{Degradation} datasets, we build upon the matched \texttt{CosmoDC2} galaxies. First, we identify SOM cells unoccupied in the \texttt{Degradation} datasets and repopulate them using corresponding galaxies from the matched \texttt{CosmoDC2} pool. Moreover, we identify galaxies that surpass the maximum redshift and magnitude limits of the \texttt{Degradation} datasets, integrating them into the \texttt{Augmentation} datasets as well, to ensure that previously excluded populations are considered.

To adaptively regulate the number of augmented galaxies, we compute \(f_\ast\), defined as the fraction of \texttt{Application} galaxies that fall into SOM cells absent from the \texttt{Degradation} dataset. We then set the target \texttt{Augmentation}-to-\texttt{Degradation} ratio to \(f_\ast \left(1  + f_\ast\right)\), thereby smoothly scaling the size of the augmented sample to balance the two datasets. In the limit of minimal spectroscopic incompleteness \(f_\ast \ll 1\), this prescription closely approximates the ideal augmentation ratio \(\frac{f_\ast}{1 - f_\ast}\). Conversely, under the most extreme degradation scenarios where \(f_\ast \to 1\), this ratio is restricted to a maximum of \(2\). This upper bound prevents the augmented sample from becoming orders of magnitude larger than its spectroscopic counterpart, thereby avoiding situations where the photometric-redshift training would be dominated by synthetic galaxies. We find that increasing the ratio beyond this cap only adds redundant augmented galaxies, raising computational costs without altering the learned colour–redshift relations or improving photometric-redshift performance.

The distribution of galaxies within the \texttt{Augmentation} datasets is also projected onto pre-trained SOM cells, as depicted in the lower left subplots of each panel in Figure~\ref{fig:3}\({-}\)\ref{fig:5}, demonstrating how data augmentation compensates for missing information within the SOM space, colour-colour diagrams, colour-magnitude diagram, respectively. Further evidence of the effectiveness of augmentation is visible in the increased presence of galaxies at \(z_\mathrm{true} \gtrsim 1.5\) (highlighted in dark-red solid lines) in Figure~\ref{fig:6}. 

As shown in Figure~\ref{fig:1}, the colour distributions of \texttt{CosmoDC2} and \texttt{OpenUniverse2024} differ markedly in \(u - g\) and \(g - r\). In contrast, Figures~\ref{fig:4} reveals that the \texttt{Augmentation} datasets exhibit substantially improved agreement with the \texttt{Application} datasets, owing to SOM‐based weighted resampling. By matching \texttt{CosmoDC2} galaxies to occupied SOM cells, we correct misalignments introduced by the distinct SED models, particularly at the bright, low-redshift end where the two catalogues overlap.

For faint, high-redshift galaxies, the differences between the SED models become more pronounced, leading to non-overlapping colour distributions, as previously shown in Figure~\ref{fig:1}. Consequently, certain galaxy populations in the \texttt{Application} datasets remain unrepresented in the \texttt{Augmentation} datasets despite weighted resampling. For example, in the LSST Y10 scenario of Figure~\ref{fig:4}, an elongated locus of galaxies with \(g - r \approx -0.5\) and \(0.0 \lesssim u - g \lesssim 3.0\) is present in the \texttt{Application} dataset but absent in the \texttt{Augmentation} dataset. Such persistent gaps can bias inferred galaxy properties and introduce systematic errors into subsequent analyses.  

\subsection{\texttt{Combination} datasets}

By merging the \texttt{Degradation} and \texttt{Augmentation} datasets, the resulting \texttt{Combination} datasets offer enhanced representativeness of galaxy broadband broadband SEDs compared to the photometric \texttt{Application} datasets, as opposed to the unaugmented \texttt{Degradation} datasets. The lower right subplots in each panel of Figure~\ref{fig:3} display the mean true redshifts of SOM cells occupied by the \texttt{Combination} datasets, demonstrating significantly enhanced coverage of the SOM cells compared to the \texttt{Degradation} datasets alone. 

Similarly, in the colour-colour and colour-magnitude diagrams presented in the lower right subplots of each panel in Figure~\ref{fig:4}-~\ref{fig:5}, combining the sparse \texttt{Degradation} datasets, which have notable deficits at fainter magnitudes and specific colour regimes, with the targeted \texttt{Augmentation} datasets results in \texttt{Combination} datasets that more effectively span the parameter space. The resultant distributions closely resemble the overall \texttt{Application} datasets, reducing incompleteness. 

Figure~\ref{fig:6} further demonstrates these gains. The under‐densities at intermediate redshifts evident in the \texttt{Degradation} samples are filled by the extended tails of the \texttt{Augmentation} distributions, yielding a well‐balanced redshift histogram in the \texttt{Combination} datasets. Here, the augmented datasets (dark-orange solid lines) more faithfully reproduce the full \texttt{Application} redshift distributions (black solid lines). 

Compared to the original implementation of data augmentation in \citet{2024ApJ...967L...6M}, which employed simplified assumptions regarding the boundaries of \(g - z\) colour, \(i-\)band magnitude and redshift, our SOM-based approach can accommodate far more complex selection functions in the multi-dimensional colour–magnitude space. 

Figure~\ref{fig:4} demonstrates the robustness of our SOM-based augmentation in colour–colour diagrams, where no regular boundary geometry is discernible. By visualising the multi-dimensional colour–magnitude space via the SOM, we can readily identify sparsely sampled regions and apply targeted data augmentation.

Similarly, as illustrated in the left panel of Figure~\ref{fig:5}, the colour–magnitude boundary of the \texttt{Degradation} dataset no longer conforms to a simple linear relation, owing to the amalgamation of multiple mock spectroscopic surveys. Nonetheless, our SOM-based augmentation adaptively compensates for gaps in spectroscopic coverage, rendering it well suited to future analyses involving complex selection functions.

However, this compensation is not entirely flawless: a noticeable deficit of galaxies around the boundaries remains. This phenomenon arises from our current definition of the \texttt{Augmentation} dataset, which augments only wholly empty SOM cells, whereas sparsely occupied cells may also require supplementary galaxies. For truly optimal performance, a comprehensive investigation of the ideal threshold for SOM cell occupancy -- potentially with adaptive thresholds for each \texttt{Degradation} dataset -- is necessary, and constitutes an avenue for future study.  

Furthermore, some SOM cells remain unpopulated or sparsely occupied even after augmentation, reflecting persistent SED model discrepancies between \texttt{CosmoDC2} and \texttt{OpenUniverse2024}. These gaps manifest most clearly as an underestimation of the high-redshift tail (\(z_\mathrm{true} \gtrsim 2.0\)), particularly in the Y10 scenario, indicating amplified mismatches for faint, distant galaxies. Moreover, variability in the spectroscopic selection functions applied during degradation translates into differing augmentation requirements across realisations, introducing an additional systematic uncertainty that we exploit to assess the impact of incomplete coverage and selection‐induced biases on photometric-redshift estimation.